\section{Evaluating the actual projected ideal trace by finite Gram data}
\label{sec:projected-ideal-gram}

The trace in MA-MO.190 contains the orthogonal projection onto the
actual diagonal ideal.  This section retains that projection.  It
gives a finite theta Gram evaluation, an explicit error certificate,
and a comparison between polynomial weights by actual nuclear maps.
It does not identify this trace with the full zero-quotient trace.

Fix the character, place set, width and one of the two regimes of
MA-MO-012.  For a polynomial exponent \(s\ge0\), write
\[
 \begin{gathered}
 w_s(x)=(e^{2ax}+e^{-2ax})(1+x^2)^{s/2},\qquad
 H_s=L^2(w_s\,dx),\qquad J_s=\overline{\mathscr I^c_\eta}^{H_s},\\
 P_s=P_{J_s},\qquad
 \mathcal A_{g,s}=\iota_s^*\mathbb T_{\eta,g}\iota_s
       \in\mathcal S_1(H_s),\qquad
 B_{g,s}=\mathcal A_{g,s}|_{J_s}.
 \end{gathered}
 \tag{MA-MO.210}
\]
The operator on the right is the original source smoothing of
MA-MO-012 on the actual completed ideal.  In particular
\(\mathcal A_{g,s}P_s\) has its range in \(J_s\), and
\(\operatorname{Tr}_{J_s}B_{g,s}=\operatorname{Tr}_{H_s}
(\mathcal A_{g,s}P_s)\).  No projection has been deleted.

\paragraph{MA-MO-014 (cross-weight comparison and a certified Gram evaluation).}
The actual ideal trace and its ordinary Fredholm determinant are
independent of \(s\ge0\).  More precisely, there are trace-class maps
\[
 B_g^{s,t}:J_s\longrightarrow J_t\quad(s,t\ge0)
 \tag{MA-MO.211}
\]
agreeing on every smooth ideal vector with the same original source
action.  Their compositions with the natural inclusions give the
endomorphisms in (MA-MO.210), and therefore
\[
 \operatorname{Tr}_{J_s}B_{g,s}=\operatorname{Tr}_{J_t}B_{g,t},
 \qquad
 \det_{J_s}(1-zB_{g,s})=\det_{J_t}(1-zB_{g,t}).
 \tag{MA-MO.212}
\]

For the two requested character families, namely nonprincipal
primitive Dirichlet characters and the characters over
\(\mathbb Q(\zeta_3)\) covered by MA-CGP-001, the same trace admits
the exact Gram-resolvent evaluation below in the metric \(w_s\).
Its finite-matrix approximation has the error bound
(MA-MO.226).  This is an evaluation of the completed-ideal trace,
with no zero-location premise added to the specified regime.

\begin{proof}
\emph{Actual maps between polynomial weights.}
Let \(k_b(t,u)\), \(b=\pm a\), be the full Fourier kernel of
\(T_{\eta,b}(g)\) in (MA-MO.187), including its compact residual.
It is Schwartz.  Use the unitary Fourier transform
\[
 V_b f(t)=(2\pi)^{-1/2}\int_{\mathbb R}e^{(b+it)x}f(x)\,dx,
 \qquad \Lambda^r=(1-\partial_t^2)^{r/2}.
\]
The map from \(H_{b,s}\) to \(H_{b,t}\) has, after these
unitary identifications, the Schwartz kernel
\[
                  \Lambda_\xi^{t/2}\Lambda_\upsilon^{-s/2}
                                           k_b(\xi,\upsilon).
 \tag{MA-MO.213}
\]
Here \(t\) is the target polynomial exponent, while
\(\xi,\upsilon\) are kernel coordinates.  Bessel multipliers of either real order
preserve Schwartz space.  The Hermite expansion thus proves
trace class between these two Hilbert spaces.

Restrict the two input components to \(J_s\).  Approximate each
input in its \(H_s\) norm by vectors of \(\mathscr I^c_\eta\).
On those vectors the two output components are the same actual
source vector, and lie in \(\mathscr I^c_\eta\).  The two
trace-class extensions converge in the respective target norms.
The closed diagonal in \(H_{a,t}\oplus H_{-a,t}\) therefore
contains the limit, which belongs to its closed ideal \(J_t\).
This proves (MA-MO.211), including its actual range assertion.
For \(t\ge s\), let \(j_{t,s}:J_t\to J_s\) be inclusion;
its norm is at most one.  Density on the common core gives
\[
 B_{g,s}=j_{t,s}B_g^{s,t},\qquad
 B_{g,t}=B_g^{s,t}j_{t,s}.
 \tag{MA-MO.214}
\]
Bounded/nuclear trace cyclicity gives the first equality in
(MA-MO.212).  The determinant equality follows first for
finite-rank approximations to \(B_g^{s,t}\) and then by
trace-norm continuity of the Fredholm determinant, exactly as
in MA-GP.75.  This argument compares actual operators; it does
not identify the orthogonal projections in different metrics.

\emph{The actual finite-band source and its exact metric.}
For the remainder of this proof fix \(s\), and write
\(H=H_s\), \(J=J_s\), \(P=P_s\), and
\(\mathcal A=\mathcal A_{g,s}\).  Use the notation of MA-CGP:
\[
 \begin{gathered}
 Q=|D_K|N\mathfrak f_\eta,\qquad
 a_n=\sum_{N\mathfrak a=n}\overline{\eta(\mathfrak a)},\qquad
 \gamma_K=2\ \text{or}\ 2\pi,\\
 \mathcal P_\Lambda=L^2((0,\Lambda),\gamma_K\,dy),\qquad
 (T_\Lambda U)(t)=
   \frac{w_K\epsilon_\eta^\theta}{\sqrt{Qt}}
                        \sum_{n\ge1}a_nU(n/(Qt)).
 \end{gathered}
 \tag{MA-MO.215}
\]
The initial domain is the actual smooth archimedean frequency
core of the required parity or radial type, with support in the
open band.  It includes the origin whenever that type permits;
no extra moment condition is imposed.  The phase and the scale
are exactly MA-CGP.6 and .16.

The target norm is now exactly
\(\int_0^\infty |F(t)|^2w_s(\log t)\,dt/t\).
The ratio of this weight to the maximum weight in MA-CGP is
between one and two.  Consequently closability and the graph
closure of \(T_\Lambda\) persist, as does its actual range
closure.  We denote the closed map by \(\overline T_\Lambda\),
its Gram operator by \(G_\Lambda=\overline T_\Lambda^*
\overline T_\Lambda\), and its range projection by \(Q_\Lambda\).
The Gram operator is formed in this exact sum metric, not copied
from an equivalent norm.

The increasing ranges have dense union in \(J\).  To see this
in the stated two-sided norm, start with any actual smooth ideal
vector.  The topological isomorphism MA-CGP.3 writes it as
\(\Theta_\eta f\).  Smooth compact frequency cutoffs preserving
the parity or radial type approximate \(f\) in Schwartz topology.
Their theta images converge in \(\mathcal E\), by the actual
isomorphism, and hence in every \(H_s\).  This proves density
without any finite-level inverse bound.  Thus
\[
 Q_\Lambda\longrightarrow P\ \text{strongly},\qquad
 \|\mathcal A Q_\Lambda-\mathcal A P\|_1\longrightarrow0.
 \tag{MA-MO.216}
\]
The second assertion follows by finite-rank approximation of
the fixed trace-class operator \(\mathcal A\).  It is not a
trace-norm assertion about the moving quotient convolutions
of MA-TN-001--002.

\emph{Normal-equation transport with all domains retained.}
Let \(\overline T_\Lambda=V_\Lambda G_\Lambda^{1/2}\) be its
polar decomposition.  For \(\varepsilon>0\), put
\[
 \begin{aligned}
 R_{\Lambda,\varepsilon}&=(G_\Lambda+\varepsilon)^{-1/2},\\
 F_{\Lambda,\varepsilon}&=\overline T_\Lambda
             R_{\Lambda,\varepsilon}
   =V_\Lambda\sqrt{G_\Lambda/(G_\Lambda+\varepsilon)}.
 \end{aligned}
 \tag{MA-MO.217}
\]
These are bounded maps on their entire source;
\(\|F_{\Lambda,\varepsilon}\|\le1\).  Their products
\(F_{\Lambda,\varepsilon}F_{\Lambda,\varepsilon}^*\)
converge strongly to \(Q_\Lambda\).  Their use does not
assert boundedness of the unregularized map \(T_\Lambda\).

Choose real compact smooth cutoffs \(0\le\chi_R\le1\),
increasing to one in the norm variable, and put
\[
 T_{\Lambda,R}=M_{\chi_R}\overline T_\Lambda,
 \qquad K_{\Lambda,R}=T_{\Lambda,R}^*\mathcal A T_{\Lambda,R}.
 \tag{MA-MO.218}
\]
The first map extends boundedly to \(\mathcal P_\Lambda\)
by the finite Poisson sum in (MA-MO.215); its values there
agree with the displayed graph-domain composition.
Consequently \(K_{\Lambda,R}\) is trace class.  Without
ever asserting that \(\overline T_\Lambda^*\mathcal A
\overline T_\Lambda\) is bounded, we have
\[
 \begin{aligned}
 R_{\Lambda,\varepsilon}K_{\Lambda,R}R_{\Lambda,\varepsilon}
 &=F_{\Lambda,\varepsilon}^*M_{\chi_R}\mathcal A
                    M_{\chi_R}F_{\Lambda,\varepsilon},\\
 \operatorname{Tr}_{J}B_{g,s}
 &=\lim_{\Lambda\to\infty}\lim_{\varepsilon\downarrow0}
       \lim_{R\to\infty}
       \operatorname{Tr}_{\mathcal P_\Lambda}
       (R_{\Lambda,\varepsilon}K_{\Lambda,R}
                                      R_{\Lambda,\varepsilon}).
 \end{aligned}
 \tag{MA-MO.219}
\]
Indeed the first line follows on the whole source because
\(R_{\Lambda,\varepsilon}\) takes its values in the graph
domain.  The error at the innermost limit is bounded by
\(\|\mathcal A-M_{\chi_R}\mathcal A M_{\chi_R}\|_1\),
uniformly in \(\Lambda,\varepsilon\).  At fixed \(\Lambda\),
strong convergence of the regularized polar maps and their
adjoints gives trace-norm convergence to
\(V_\Lambda^*\mathcal A V_\Lambda\).  Cyclicity gives its
trace as \(\operatorname{Tr}(\mathcal A Q_\Lambda)\), and
(MA-MO.216) proves the outer limit.  No unrestricted joint
choice of \(\varepsilon\) and \(\Lambda\) is implicit.

\emph{An explicit finite local kernel.}
Write \(\mathfrak a(t,z)\) for the kernel of \(\mathcal A\)
relative to \(dz/z\).  It is smooth on compact positive
rectangles.  Explicitly, if \(\ell_b(x,y)\) is the inverse
double Fourier transform of \(k_b\), then
\[
 \begin{aligned}
 \mathfrak a(e^x,e^y)
   &=\sum_{b=\pm a}\theta_b(x)e^{-b(x-y)}\ell_b(x,y),\\
 \theta_b(x)&=\frac{e^{2bx}}{e^{2ax}+e^{-2ax}},\qquad
 \ell_b(x,y)=\frac1{2\pi}\iint
                 e^{-itx}k_b(t,u)e^{iuy}\,dt\,du .
 \end{aligned}
 \tag{MA-MO.220}
\]
The adjoint of the inclusion \(H_s\to H_{b,s}\) is
precisely multiplication by \(\theta_b\).  This proves
(MA-MO.220) in the exact metric, including its normalization.

Substitute (MA-MO.215) twice into (MA-MO.218), and set
\(y=m/(Qt)\), \(v=n/(Qz)\).  The resulting action, with
integration measure \(dv\), is
\[
 \begin{aligned}
 (K_{\Lambda,R}U)(y)
  ={}&\frac{w_K^2}{\gamma_K}
       \sum_{m,n\ge1}\frac{\overline{a_m}a_n}{\sqrt{mn}}
       \int_0^\Lambda
       \frac{w_s(\log(m/(Qy)))}{\sqrt{yv}}\\
     &\quad\cdot\chi_R(m/(Qy))\chi_R(n/(Qv))
       \mathfrak a(m/(Qy),n/(Qv))U(v)\,dv .
 \end{aligned}
 \tag{MA-MO.221}
\]
If the cutoff is supported in \([r_-,r_+]\subset(0,\infty)\),
only \(m,n\le Q\Lambda r_+\) occur, and every nonzero
term is away from \(y=0\) and \(v=0\).  This is a finite
smooth-kernel formula on the actual auxiliary band space.
Its trace-class property also follows directly by extending
each kernel smoothly to a larger square and using an
absolutely summable double Fourier series.  The scalar theta
phase cancels; the primitive scale \(Q\), the conjugate
coefficients, and the exact target weight do not.

For reference the localized positive Gram form uses the
finite rational-dilation operator
\[
 (G_{\Lambda,R}U)(y)=\frac{w_K^2}{\gamma_K}
       \sum_{m,n\ge1}\frac{\overline{a_m}a_n}{m}
       \chi_R(m/(Qy))^2w_s(\log(m/(Qy)))U(ny/m),
 \tag{MA-MO.222}
\]
with extension by zero off the band.  Its forms increase
on the graph core to \(\|\overline T_\Lambda U\|_H^2\).
This supplies the actual normal equation in (MA-MO.217);
an uncut double sum is not interchanged with an integral.

\emph{A finite-matrix evaluation with a stopping certificate.}
Choose a countable sequence of actual smooth frequency tests,
each lying in some finite band, whose theta images \(v_j\)
have dense span in \(J\).  Such a sequence follows from the
frequency approximation just proved and separability of the
Schwartz source; delete linear dependencies.  Let
\(V_N:\mathbb C^N\to H\) have columns \(v_1,\ldots,v_N\),
and define
\[
 \begin{gathered}
 C_N=V_N^*V_N>0,\qquad D_N=V_N^*\mathcal A V_N,\qquad
 P_N=V_NC_N^{-1}V_N^*,\\
 \tau_N=\operatorname{tr}(C_N^{-1}D_N)
                   =\operatorname{Tr}_H(\mathcal A P_N).
 \end{gathered}
 \tag{MA-MO.223}
\]
Every matrix entry is an actual theta pairing.  It can be
computed by compact truncation using (MA-MO.221)--(MA-MO.222),
with the Schwartz tails of these finitely many smooth vectors
retained in the truncation error.

The complementary data are equally explicit.  For any finite
set of retained zero jets from MA-WP.3, take their Riesz vectors
\[
 h_{\rho,j}(x)=\frac{x^j e^{\overline{(\rho-1/2)}x}}{w_s(x)}.
 \tag{MA-MO.224}
\]
The inner product is linear in its first variable, so
\(\langle f,h_{\rho,j}\rangle=(\mathcal Mf)^{(j)}(\rho-1/2)\).
These vectors belong to \(J^\perp\) in the exact metric.
Their Gram projection \(E_M=W_M(W_M^*W_M)^{-1}W_M^*\)
is defined by their column map.  Distinct retained jets are
linearly independent, since exponential polynomials are
linearly independent.  If there are no retained jets use
\(E_M=0\).  The full exhaustive family spans \(J^\perp\)
by MA-WP.4, which applies to the sum norm as stated there.
In particular \(P_NE_M=0\).

Here is a completely specified nuclear approximation of
\(\mathcal A\).  Put \(W_{b,s}=\Lambda^{s/2}V_b\), let
\(e_m\) be the normalized Hermite basis of \(L^2(\mathbb R)\),
and let \(c^b_{mn}\) be the Hermite matrix of the Schwartz
kernel \(\Lambda_t^{s/2}\Lambda_u^{-s/2}k_b\).
If \(i_b:H\to H_{b,s}\) is inclusion, put
\(u^b_m=i_b^*W_{b,s}^{-1}e_m\).  Then
\[
 \begin{gathered}
 \mathcal A=\sum_{b=\pm a}\sum_{m,n\ge0}
      c^b_{mn}\,u^b_m\otimes(u^b_n)^*,\qquad \|u^b_m\|\le1,\\
 \delta_{N,M}(u)=\|u\|^2-\|P_Nu\|^2-\|E_Mu\|^2\ge0 .
 \end{gathered}
 \tag{MA-MO.225}
\]
All sums converge in trace norm.  If
\(e_L=\sum_b\sum_{\max(m,n)>L}|c^b_{mn}|\), the error bound is
\[
 \left|\operatorname{Tr}_{J}B_{g,s}-\tau_N\right|
 \le e_L+
 \sum_{b=\pm a}\sum_{m,n\le L}|c^b_{mn}|
       \sqrt{\delta_{N,M}(u^b_m)\delta_{N,M}(u^b_n)} .
 \tag{MA-MO.226}
\]
To prove it, \(D=P-P_N\) is an orthogonal projection and
\(0\le D\le1-P_N-E_M\).  Hence
\( |\langle Du,v\rangle|
 \le\sqrt{\delta_{N,M}(u)\delta_{N,M}(v)}\).
Apply this to the finite part of (MA-MO.225); the remaining
trace is bounded by its nuclear tail.  Every quantity on
the right uses finite Gram matrices and explicitly specified
kernel coefficients.  An incomplete finite set of certified
zero jets still gives a valid upper bound; it is not assumed
to list all zeros.

The nuclear tail has a concrete seminorm majorant.  With
\(\mathcal H=-\partial_t^2+t^2\), for any integer \(q>1\)
let
\[
 C_{b,q}=\| (\mathcal H_t+1)^q(\mathcal H_u+1)^q
           (\Lambda_t^{s/2}\Lambda_u^{-s/2}k_b)\|_{L^2(\mathbb R^2)}.
\]
Integration against Hermite eigenfunctions gives
\[
 \begin{aligned}
 |c^b_{mn}|&\le C_{b,q}(2m+2)^{-q}(2n+2)^{-q},\\
 e_L&\le\sum_b C_{b,q}\left[
       \left(\sum_{m\ge0}(2m+2)^{-q}\right)^2
       -\left(\sum_{m=0}^L(2m+2)^{-q}\right)^2\right].
 \end{aligned}
 \tag{MA-MO.227}
\]
Thus the bound does not hide an unknown projector norm.
For fixed \(L\), \(P_N\to P\) and \(E_M\to1-P\) strongly
make its finite part tend to zero.  Then \(L\to\infty\)
makes its tail zero.  This proves a convergent finite-matrix
evaluation of the specified trace.
\end{proof}

\paragraph{What is and is not evaluated.}
Equations (MA-MO.219) and (MA-MO.223)--(MA-MO.227) evaluate
\(\operatorname{Tr}_{J_s}B_{g,s}\), retaining the theta Gram
normalization and the actual ideal projection.  The bound is
an a posteriori analytic certificate when its finite pairings,
certified zero jets and kernel-seminorm bounds are supplied.
The full-measure generic-width theorem by itself does not
provide an effective zero-distance constant for an arbitrarily
chosen decimal width, and no uniform numerical convergence
rate is claimed here.

The source inputs are MA-CGP-001--003, especially .3 and
.15--.19; MA-GP-005/008 for the closed normal equation and
bounded regularization; MA-WP-001 for the exact complementary
jet span; and the actual kernels and source agreement of
MA-MO-012.  The finite matrices use the inherited two-sided
metric throughout.  They are Hilbert Gram matrices, not
categorical identifications of different completions.  Their
trace limits use the actual nuclear pairing of the inherited
cote~1 framework.  The retained jet spaces are finite-length
primary modules, but their algebraic classes alone do not
determine this Gram correction.

No equality with \(\operatorname{Tr}C_g\) on the full smooth
zero quotient follows from this evaluation.  Such an equality
requires an additional actual comparison or return map.
In particular no part of (MA-MO.212) or (MA-MO.226) removes
the omitted-divisor condition of MA-FT-002.
